\documentclass[letterpaper]{article}
\usepackage[utf8]{inputenc}
\usepackage[spanish,es-nodecimaldot]{babel}
\usepackage[T1]{fontenc}
\usepackage{mathptmx}
% Fuente monoespaciada sin ligaduras: evita que "--opcion" se imprima como una sola raya
\renewcommand{\ttdefault}{pcr}
\usepackage{microtype}
\DisableLigatures[-]{encoding = *, family = tt*}
\usepackage[top=2cm, bottom=2cm, left=2.5cm, right=2.5cm]{geometry}
\usepackage{titlesec}
\usepackage{parskip}
\usepackage{enumitem}
\usepackage{graphicx}
\usepackage{fancyhdr}
\usepackage{hyperref}
\usepackage{amsmath}
\usepackage{array}
\usepackage{tabularx}
\usepackage{caption}
\usepackage{float}
\usepackage{setspace}
\usepackage{multirow}
\usepackage{colortbl}
\usepackage{diagbox}
\usepackage{pifont}
\usepackage{makecell}
\usepackage{listings}

% Cargar TikZ y librerías una sola vez
\usepackage{tikz}
\usetikzlibrary{shapes.geometric, arrows.meta, positioning, fit, backgrounds, calc, shadows}

% Paleta de colores (misma del protocolo)
\usepackage[table]{xcolor}
\definecolor{colorDatos}{RGB}{173, 216, 230}
\definecolor{colorModulos}{RGB}{144, 238, 144}
\definecolor{colorProcesos}{RGB}{255, 239, 178}
\definecolor{colorTexto}{RGB}{30, 30, 30}
\definecolor{colorPropuesta}{RGB}{230, 240, 255}
\definecolor{colorOtros}{RGB}{245, 245, 245}
\definecolor{semNormal}{RGB}{200, 230, 201}
\definecolor{semPrecaucion}{RGB}{255, 243, 196}
\definecolor{semCritico}{RGB}{255, 205, 210}

\newcommand{\cmark}{\ding{51}}
\newcommand{\xmark}{\ding{55}}

\addto\captionsspanish{\renewcommand{\tablename}{Tabla}}
\addto\captionsspanish{\renewcommand{\listtablename}{Índice de tablas}}

% ===== Utilidades de redacción =====
% Comandos de terminal (admite espacios y guiones; los espacios permiten salto de línea)
\newcommand{\cmd}[1]{\texttt{\detokenize{#1}}}
% Rutas y nombres de archivo (admite guion bajo y salto de línea en las diagonales)
\newcommand{\ruta}[1]{\path{#1}}
\emergencystretch=3em

% Bloques de código / terminal
\lstset{
  basicstyle=\ttfamily\footnotesize,
  backgroundcolor=\color{colorOtros},
  frame=single,
  rulecolor=\color{gray!60},
  breaklines=true,
  columns=fullflexible,
  keepspaces=true,
  showstringspaces=false,
  xleftmargin=4pt, xrightmargin=4pt,
  aboveskip=6pt, belowskip=6pt
}

% Estilo de encabezados (equivalente al del protocolo)
\titleformat{\section}{\fontsize{12}{14}\bfseries\selectfont}{\thesection.}{0.6em}{}
\titleformat{\subsection}{\fontsize{11}{13}\bfseries\selectfont}{\thesubsection}{0.6em}{}
\titlespacing*{\section}{0pt}{14pt}{6pt}
\titlespacing*{\subsection}{0pt}{10pt}{4pt}

\hypersetup{colorlinks=true, linkcolor=blue!50!black, urlcolor=blue!50!black, citecolor=blue!50!black}

% Encabezado y pie de página
\pagestyle{fancy}
\fancyhf{}
\fancyhead[L]{\fontsize{8}{10}\selectfont\itshape Manual técnico y de usuario · TT No. 2027-A190}
\fancyhead[R]{\fontsize{8}{10}\selectfont\itshape ESCOM-IPN}
\fancyfoot[C]{\fontsize{9}{11}\selectfont\thepage}
\renewcommand{\headrulewidth}{0.3pt}

% Ayudas para tablas con el estilo del protocolo
\newcommand{\encabezadotabla}{\rowcolor{red!20}}
\newcommand{\filatabla}{\rowcolor{colorOtros}}

% ============================================================
\begin{document}

\begin{center}
    {\fontsize{12}{15}\bfseries\selectfont
    Manual técnico y de usuario del prototipo basado en la aplicación de
    aprendizaje profundo para estimar el tiempo de vida útil restante en motores
    turbofan de aeronave comercial Boeing~777}\\[6pt]
    {\fontsize{12}{15}\bfseries\itshape\selectfont
    Trabajo Terminal No.~2027-A190}\\[6pt]
    {\fontsize{10}{13}\itshape\selectfont
    Alumnos: *Rojas Vasquez Evelyne$^{1}$, Onofre Rosales Rodrigo$^{2}$ \\
    Directores: Morales Guitron Sandra Luz$^{1}$, López Ruiz Gabriela de Jesús$^{2}$ \\
    *e-mail: erojasv2200@alumno.ipn.mx}
\end{center}

\vspace{4pt}
\begin{table}[H]
\centering
\fontsize{9}{11}\selectfont
\renewcommand{\arraystretch}{1.4}
\begin{tabularx}{\textwidth}{|>{\bfseries}l|X|}
    \hline
    \encabezadotabla \multicolumn{2}{|c|}{\textbf{Ficha del documento}} \\
    \hline
    \filatabla Documento & Manual técnico y de usuario (entregable de la fase de Documentación del cronograma) \\
    \hline
    \filatabla Versión & 1.0 — octubre de 2026 \\
    \hline
    \filatabla Ubicación en el repositorio & \ruta{docs/manual_usuario/manual_usuario.tex} (fuente) y \ruta{docs/manual_usuario/manual_usuario.pdf} (compilado) \\
    \hline
    \filatabla Corresponde a & Código del repositorio versión 1.0.0 (\ruta{src/__init__.py}) \\
    \hline
    \filatabla Dirigido a & Operadores del departamento de mantenimiento aeronáutico (uso del tablero) y personal técnico (instalación, entrenamiento y evaluación) \\
    \hline
\end{tabularx}
\end{table}

\newpage
\tableofcontents
\newpage
% ============================================================
\section{Introducción}

Este documento explica cómo instalar, ejecutar y utilizar el \textbf{prototipo}
de sistema de mantenimiento predictivo desarrollado en el Trabajo Terminal
No.~2027-A190. El prototipo \textbf{aplica} una arquitectura de red neuronal
recurrente LSTM reportada en la literatura~\cite{zheng2017lstm} para estimar la
vida útil restante (\textit{Remaining Useful Life}, RUL) de motores turbofan a
partir de lecturas multivariadas de sensores, y presenta los resultados en un
tablero de visualización con un indicador de alerta de falla.

Siguiendo la terminología del protocolo, el término \textbf{prototipo} designa
el sistema integrado (módulos de software, flujo de datos y tablero), mientras
que el término \textbf{modelo} designa la red LSTM que funciona como núcleo de
inferencia.

\subsection{Alcance y advertencia sobre los datos}

El caso de aplicación es el motor turbofan de aeronaves comerciales como el
Boeing~777. Sin embargo, el prototipo se entrena y valida con los conjuntos de
datos públicos C-MAPSS de la NASA~\cite{saxena2008damage}, que son resultado de
una \textbf{simulación} de un motor turbofan genérico~\cite{frederick2007cmapss}.
Por lo tanto:

\begin{itemize}
    \item Los datos \textbf{no} son telemetría real de motores Boeing~777.
    \item Los canales disponibles son temperaturas, presiones, velocidades,
          razones y flujos (Tabla~\ref{tab:sensores}); \textbf{no} existe un canal
          de vibración.
    \item Los resultados son de carácter académico y \textbf{no} sustituyen los
          procedimientos de mantenimiento ni de aeronavegabilidad vigentes.
\end{itemize}

\subsection{Qué encontrará en este manual}

\begin{table}[H]
\centering
\fontsize{9}{11}\selectfont
\caption{Guía de lectura según el perfil del usuario.}
\label{tab:guia}
\renewcommand{\arraystretch}{1.5}
\begin{tabularx}{\textwidth}{|>{\raggedright\arraybackslash}p{4.3cm}|X|>{\raggedright\arraybackslash}p{3cm}|}
    \hline
    \encabezadotabla \textbf{Perfil} & \textbf{Qué necesita hacer} & \textbf{Secciones} \\
    \hline
    \filatabla Operador de mantenimiento & Consultar el estado de un motor o de la flota e interpretar las alertas & \ref{sec:tablero}, \ref{sec:limitaciones} \\
    \hline
    \filatabla Personal técnico / desarrollador & Instalar, obtener los datos, entrenar, evaluar y mantener el sistema & \ref{sec:requisitos}--\ref{sec:entrenamiento}, \ref{sec:estructura}, \ref{sec:problemas} \\
    \hline
    \filatabla Evaluador (sinodal) & Reproducir los resultados y verificar la metodología & \ref{sec:descripcion}, \ref{sec:entrenamiento}, \ref{sec:estructura} \\
    \hline
\end{tabularx}
\end{table}

% ============================================================
\section{Descripción general del sistema}
\label{sec:descripcion}

El prototipo se organiza en cuatro módulos que operan en secuencia, tal como se
describe en el protocolo. La Figura~\ref{fig:arquitectura} muestra el flujo
completo: a la izquierda, la \textbf{entrada} (datos crudos de sensores del
motor); a la derecha, la \textbf{salida} (valor estimado de RUL en ciclos de
operación y alerta visual en el tablero).

\begin{figure}[H]
\centering
\resizebox{\textwidth}{!}{%
\begin{tikzpicture}[
    io/.style={
      rectangle, rounded corners=3pt,
      minimum width=3.6cm, minimum height=1.5cm,
      text width=3.6cm, align=center,
      draw=black!70, fill=gray!15,
      font=\fontsize{9}{11}\selectfont\bfseries
    },
    modulo/.style={
      rectangle, rounded corners=5pt,
      minimum width=3.6cm, minimum height=2.0cm,
      text width=3.6cm, align=center,
      draw=gray!60, fill=#1,
      font=\fontsize{9}{11}\selectfont,
      drop shadow={shadow xshift=1.5pt, shadow yshift=-1.5pt, fill=gray!30}
    },
    etiqueta/.style={
      font=\fontsize{7.5}{9}\selectfont\itshape,
      text width=2.1cm, align=center
    },
    flecha/.style={-Stealth, thick, draw=gray!60},
    caja/.style={
      rectangle, rounded corners=7pt,
      draw=gray!40, dashed, thick,
      fill=gray!4,
      inner xsep=10pt, inner ysep=12pt
    }
  ]
  % Fila 1
  \node[io] (entrada) at (0,0) {%
    \textbf{ENTRADA}\\[3pt]
    Lecturas de sensores\\del motor turbofan\\(dataset C-MAPSS)
  };
  \node[modulo=blue!12] (adq) at (5.8,0) {%
    \textbf{Módulo 1}\\[4pt]
    Adquisición y\\limpieza de datos\\[3pt]
    {\fontsize{7}{8}\selectfont\itshape(eliminación de señales\\constantes, normalización)}
  };
  \node[modulo=green!14] (ventanas) at (11.6,0) {%
    \textbf{Módulo 2}\\[4pt]
    Construcción de\\ventanas temporales\\[3pt]
    {\fontsize{7}{8}\selectfont\itshape(segmentación de series\\de tiempo por motor)}
  };
  % Fila 2
  \node[modulo=orange!18] (lstm) at (0,-5.2) {%
    \textbf{Módulo 3}\\[4pt]
    Modelo de\\estimación LSTM\\[3pt]
    {\fontsize{7}{8}\selectfont\itshape(red neuronal recurrente\\para predecir el RUL)}
  };
  \node[modulo=red!10] (salida_mod) at (5.8,-5.2) {%
    \textbf{Módulo 4}\\[4pt]
    Generación de\\resultados y alertas\\[3pt]
    {\fontsize{7}{8}\selectfont\itshape(cálculo de métricas\\RMSE, MAE y Score)}
  };
  \node[io] (salida) at (11.6,-5.2) {%
    \textbf{SALIDA}\\[3pt]
    Valor estimado de\\vida útil restante\\(RUL en ciclos)
  };
  % Tablero
  \node[io, fill=purple!12] (dash) at (11.6,-9.2) {%
    Tablero de visualización\\[2pt]
    {\fontsize{7.5}{9}\selectfont\itshape(gráficas + indicador\\de alerta de falla)}
  };

  % Flechas fila 1
  \draw[flecha] (entrada) -- (adq)
      node[etiqueta, midway, above=2pt]{datos crudos\\de sensores};
  \draw[flecha] (adq) -- (ventanas)
      node[etiqueta, midway, above=2pt]{series limpias\\y normalizadas};
  % Salto de fila: Módulo 2 -> Módulo 3
  \draw[flecha] (ventanas.south) -- ++(0,-1.5) -| (lstm.north);
  \node[etiqueta, text width=4cm, fill=white, inner sep=1pt] at (5.8,-2.2) {ventanas de tamaño fijo};
  % Flechas fila 2
  \draw[flecha] (lstm) -- (salida_mod)
      node[etiqueta, midway, above=2pt]{RUL estimado\\por motor};
  \draw[flecha] (salida_mod) -- (salida)
      node[etiqueta, midway, above=2pt]{valor numérico\\de RUL};
  \draw[flecha] (salida.south) -- (dash.north)
      node[midway, right=3pt, font=\fontsize{7.5}{9}\selectfont\itshape]{visualización};
  % Retroalimentación
  \draw[flecha, dashed, draw=gray!55] (salida_mod.south) -- ++(0,-1.0) -| (lstm.south);
  \node[font=\fontsize{7}{8.5}\selectfont\itshape, align=center, anchor=north] at (2.9,-7.25)
      {retroalimentación de métricas\\(ajuste de hiperparámetros)};

  \begin{pgfonlayer}{background}
    \node[caja, fit=(adq)(ventanas),
          label={[font=\fontsize{8.5}{10}\selectfont\bfseries,
                  text=blue!50!black]above:Capa de preprocesamiento}] {};
    \node[caja, fit=(lstm)(salida_mod),
          label={[font=\fontsize{8.5}{10}\selectfont\bfseries,
                  text=orange!60!black]above:Capa de inferencia y evaluación}] {};
  \end{pgfonlayer}
\end{tikzpicture}%
}
\caption{Arquitectura modular del sistema de estimación de vida útil restante
(RUL). La \textbf{entrada} son los datos de sensores del motor; la
\textbf{salida} es el número de ciclos de operación restantes estimados por el
modelo. El flujo se lee de izquierda a derecha en cada fila y continúa en la
fila inferior.}
\label{fig:arquitectura}
\end{figure}

La Tabla~\ref{tab:modulos} indica en qué archivos del repositorio se implementa
cada módulo.

\begin{table}[H]
\centering
\fontsize{9}{11}\selectfont
\caption{Correspondencia entre los módulos del protocolo y el código fuente.}
\label{tab:modulos}
\renewcommand{\arraystretch}{1.5}
\begin{tabularx}{\textwidth}{|>{\raggedright\arraybackslash}p{3.6cm}|>{\raggedright\arraybackslash}p{4.6cm}|X|}
    \hline
    \encabezadotabla \textbf{Módulo} & \textbf{Archivos} & \textbf{Responsabilidad} \\
    \hline
    \filatabla 1. Adquisición y limpieza & \ruta{src/ingestion.py}, \ruta{src/preprocessing.py} & Descarga y lectura validada de C-MAPSS; eliminación de sensores constantes; normalización; etiquetas de RUL \\
    \hline
    \filatabla 2. Ventanas temporales & \ruta{src/preprocessing.py} & Ventanas deslizantes por motor; partición entrenamiento/validación por motor \\
    \hline
    \filatabla 3. Modelo LSTM & \ruta{src/model.py}, \ruta{src/train.py} & Arquitectura LSTM; optimizador Adam con reducción adaptativa de la tasa de aprendizaje; entrenamiento \\
    \hline
    \filatabla 4. Resultados y alertas & \ruta{src/evaluation.py}, \ruta{src/inference.py}, \ruta{dashboard/app.py} & Métricas RMSE, MAE y Score; gráficas; curva de RUL; alertas; tablero \\
    \hline
\end{tabularx}
\end{table}

\subsection{Metodología aplicada}

El desarrollo sigue el estándar CRISP-DM~\cite{wirth2000crisp}. Las seis fases se
agrupan en cuatro etapas ejecutables que se reflejan en el uso del sistema
(Tabla~\ref{tab:crisp}).

\begin{table}[H]
\centering
\fontsize{9}{11}\selectfont
\caption{Etapas del proyecto, fases CRISP-DM y comando asociado.}
\label{tab:crisp}
\renewcommand{\arraystretch}{1.5}
\begin{tabularx}{\textwidth}{|>{\raggedright\arraybackslash}p{4.2cm}|>{\centering\arraybackslash}p{2.4cm}|X|}
    \hline
    \encabezadotabla \textbf{Etapa} & \textbf{Fases CRISP-DM} & \textbf{Dónde se ejecuta} \\
    \hline
    \filatabla 1. Recopilación y preparación de datos & 1, 2 y 3 & \cmd{make data}; cuaderno \ruta{notebooks/01_EDA.ipynb} \\
    \hline
    \filatabla 2. Adaptación y entrenamiento del modelo & 4 & \cmd{make train}, \cmd{make grid} \\
    \hline
    \filatabla 3. Evaluación y validación & 5 & Evaluación automática al final de \cmd{make train}; \cmd{make evaluate} \\
    \hline
    \filatabla 4. Despliegue: tablero de visualización & 6 & \cmd{make dashboard} \\
    \hline
\end{tabularx}
\end{table}

% ============================================================
\section{Requisitos}
\label{sec:requisitos}

\begin{table}[H]
\centering
\fontsize{9}{11}\selectfont
\caption{Requisitos de hardware y software.}
\label{tab:requisitos}
\renewcommand{\arraystretch}{1.5}
\begin{tabularx}{\textwidth}{|>{\raggedright\arraybackslash}p{3.8cm}|X|}
    \hline
    \encabezadotabla \textbf{Componente} & \textbf{Requisito} \\
    \hline
    \filatabla Sistema operativo & Windows 10/11, macOS o Linux \\
    \hline
    \filatabla Python & Versión 3.10 a 3.12 (con \cmd{pip} y \cmd{venv}) \\
    \hline
    \filatabla Memoria RAM & 8~GB recomendados (el modelo es pequeño; el entrenamiento es viable sin GPU) \\
    \hline
    \filatabla Disco & Aproximadamente 3~GB para dependencias (TensorFlow) y datos \\
    \hline
    \filatabla Bibliotecas principales & TensorFlow $\geq$ 2.16 (Keras~3), NumPy, Pandas, Scikit-Learn, Matplotlib, Plotly y Streamlit $\geq$ 1.50; la lista completa está en \ruta{requirements.txt} \\
    \hline
    \filatabla Herramientas opcionales & \cmd{make} (en Windows puede sustituirse, ver Tabla~\ref{tab:equivalencias}); Git; \LaTeX{} (\cmd{latexmk}) para compilar este manual \\
    \hline
    \filatabla Conexión a Internet & Necesaria para instalar dependencias y descargar C-MAPSS \\
    \hline
\end{tabularx}
\end{table}

% ============================================================
\section{Instalación paso a paso}
\label{sec:instalacion}

\begin{figure}[H]
\centering
\resizebox{\textwidth}{!}{%
\begin{tikzpicture}[
    paso/.style={
      rectangle, rounded corners=5pt,
      minimum width=2.9cm, minimum height=1.7cm,
      text width=2.9cm, align=center,
      draw=gray!60, fill=#1,
      font=\fontsize{8.5}{10.5}\selectfont,
      drop shadow
    },
    fl/.style={-Stealth, thick, draw=gray!60}
  ]
  \node[paso=blue!9] (p1) {\textbf{Paso 1}\\[2pt]Instalar\\dependencias\\[2pt]{\fontsize{7}{9}\selectfont\ttfamily make setup}};
  \node[paso=green!12, right=0.9cm of p1] (p2) {\textbf{Paso 2}\\[2pt]Obtener\\los datos C-MAPSS\\[2pt]{\fontsize{7}{9}\selectfont\ttfamily make data}};
  \node[paso=orange!16, right=0.9cm of p2] (p3) {\textbf{Paso 3}\\[2pt]Entrenar y\\evaluar el modelo\\[2pt]{\fontsize{7}{9}\selectfont\ttfamily make train}};
  \node[paso=yellow!22, right=0.9cm of p3] (p4) {\textbf{Paso 4}\\[2pt]Revisar\\las métricas\\[2pt]{\fontsize{7}{9}\selectfont\ttfamily metrics.json}};
  \node[paso=purple!12, right=0.9cm of p4] (p5) {\textbf{Paso 5}\\[2pt]Abrir el\\tablero\\[2pt]{\fontsize{7}{9}\selectfont\ttfamily make dashboard}};
  \draw[fl] (p1)--(p2);
  \draw[fl] (p2)--(p3);
  \draw[fl] (p3)--(p4);
  \draw[fl] (p4)--(p5);
\end{tikzpicture}%
}
\caption{Flujo de uso del prototipo, desde la instalación hasta el tablero.}
\label{fig:flujo}
\end{figure}

\subsection{Paso 1. Obtener el código}

\begin{lstlisting}
git clone <URL_DEL_REPOSITORIO> rul-turbofan-tt
cd rul-turbofan-tt
\end{lstlisting}

Todos los comandos de este manual se ejecutan desde la \textbf{raíz del
repositorio} (la carpeta que contiene \ruta{Makefile} y \ruta{README.md}).

\subsection{Paso 2. Crear y activar un entorno virtual}

\textbf{Linux o macOS:}
\begin{lstlisting}
python -m venv .venv
source .venv/bin/activate
\end{lstlisting}

\textbf{Windows (PowerShell):}
\begin{lstlisting}
python -m venv .venv
.venv\Scripts\activate
\end{lstlisting}

El indicador del entorno (\cmd{(.venv)}) debe aparecer al inicio de la línea de
comandos.

\subsection{Paso 3. Instalar las dependencias}

\begin{lstlisting}
make setup
\end{lstlisting}

Este comando ejecuta \cmd{python -m pip install -r requirements-dev.txt}, que
incluye las dependencias de ejecución, \textit{pytest}, \textit{ruff} y
JupyterLab. La instalación de TensorFlow puede tardar varios minutos.

\subsection{Paso 4. Verificar la instalación}

\begin{lstlisting}
make test
\end{lstlisting}

Deben ejecutarse 24 pruebas unitarias y de extremo a extremo (usan datos
sintéticos, por lo que no requieren descargar C-MAPSS). La salida correcta
termina con una línea similar a \cmd{24 passed}.

\begin{table}[H]
\centering
\fontsize{9}{11}\selectfont
\caption{Equivalencia entre comandos \cmd{make} y comandos directos (útil en Windows sin \cmd{make}).}
\label{tab:equivalencias}
\renewcommand{\arraystretch}{1.5}
\begin{tabularx}{\textwidth}{|>{\raggedright\arraybackslash}p{4.2cm}|X|}
    \hline
    \encabezadotabla \textbf{Con make} & \textbf{Comando directo} \\
    \hline
    \filatabla \cmd{make setup} & \cmd{python -m pip install -r requirements-dev.txt} \\
    \hline
    \filatabla \cmd{make data} & \cmd{python -m src.ingestion --download} \\
    \hline
    \filatabla \cmd{make test} & \cmd{python -m pytest -q} \\
    \hline
    \filatabla \cmd{make train SUBSET=FD001 WINDOW=30} & \cmd{python -m src.train --subset FD001 --window 30 --epochs 100} \\
    \hline
    \filatabla \cmd{make grid SUBSET=FD001} & \cmd{python -m src.grid_search --subset FD001 --epochs 100} \\
    \hline
    \filatabla \cmd{make evaluate SUBSET=FD001} & \cmd{python -m src.evaluation --subset FD001} \\
    \hline
    \filatabla \cmd{make dashboard} & \cmd{streamlit run dashboard/app.py} \\
    \hline
    \filatabla \cmd{make manual} & \cmd{cd docs/manual_usuario} y luego \cmd{latexmk -pdf manual_usuario.tex} \\
    \hline
\end{tabularx}
\end{table}

% ============================================================
\section{Obtención y descripción de los datos}
\label{sec:datos}

\subsection{Descarga}

\begin{lstlisting}
make data
\end{lstlisting}

El comando descarga el archivo comprimido de C-MAPSS, lo extrae en
\ruta{data/raw/} y verifica que existan los 12 archivos esperados. Si la
descarga falla (la ubicación del recurso en los servidores de la NASA ha
cambiado en el pasado), hay dos alternativas:

\begin{enumerate}[leftmargin=*, label=\arabic*.]
    \item Descargar manualmente el conjunto desde el repositorio oficial de datos
          de pronóstico de la NASA~\cite{saxena2008damage} y copiar los archivos
          \ruta{train_FD00X.txt}, \ruta{test_FD00X.txt} y \ruta{RUL_FD00X.txt}
          (con X = 1, 2, 3, 4) en \ruta{data/raw/}.
    \item Indicar la dirección vigente: \cmd{python -m src.ingestion --download --url <URL>}.
\end{enumerate}

Para comprobar la carga, ejecute \cmd{python -m src.ingestion}: imprime el
número de motores y de filas de cada subconjunto.

\subsection{Subconjuntos y archivos}

\begin{table}[H]
\centering
\fontsize{9}{11}\selectfont
\caption{Subconjuntos de C-MAPSS~\cite{saxena2008damage}.}
\label{tab:subconjuntos}
\renewcommand{\arraystretch}{1.5}
\begin{tabularx}{\textwidth}{|c|c|c|c|c|}
    \hline
    \encabezadotabla \textbf{Subconjunto} & \textbf{Regímenes de operación} & \textbf{Modos de falla} & \textbf{Motores de entrenamiento} & \textbf{Motores de prueba} \\
    \hline
    \filatabla FD001 & 1 & 1 & 100 & 100 \\
    \hline
    \filatabla FD002 & 6 & 1 & 260 & 259 \\
    \hline
    \filatabla FD003 & 1 & 2 & 100 & 100 \\
    \hline
    \filatabla FD004 & 6 & 2 & 249 & 248 \\
    \hline
\end{tabularx}
\end{table}

Cada subconjunto consta de tres archivos de texto:

\begin{itemize}
    \item \ruta{train_FD00X.txt}: trayectorias completas de cada motor, desde el
          inicio de la operación hasta la falla.
    \item \ruta{test_FD00X.txt}: trayectorias truncadas antes de la falla.
    \item \ruta{RUL_FD00X.txt}: RUL real del último ciclo observado de cada
          motor de prueba (una línea por motor).
\end{itemize}

Cada línea de \ruta{train_} y \ruta{test_} contiene 26 columnas separadas por
espacios: identificador del motor (\cmd{unit}), ciclo (\cmd{cycle}), tres ajustes
operacionales (\cmd{op1}--\cmd{op3}) y 21 canales de sensores (\cmd{s1}--\cmd{s21}).

\begin{table}[H]
\centering
\fontsize{8.5}{10.5}\selectfont
\caption{Canales de sensores de C-MAPSS~\cite{frederick2007cmapss,saxena2008damage}.}
\label{tab:sensores}
\renewcommand{\arraystretch}{1.25}
\begin{tabularx}{\textwidth}{|c|l|X|}
    \hline
    \encabezadotabla \textbf{Canal} & \textbf{Símbolo} & \textbf{Descripción} \\
    \hline
    \filatabla s1 & T2 & Temperatura total a la entrada del fan \\ \hline
    \filatabla s2 & T24 & Temperatura total a la salida del compresor de baja (LPC) \\ \hline
    \filatabla s3 & T30 & Temperatura total a la salida del compresor de alta (HPC) \\ \hline
    \filatabla s4 & T50 & Temperatura total a la salida de la turbina de baja (LPT) \\ \hline
    \filatabla s5 & P2 & Presión a la entrada del fan \\ \hline
    \filatabla s6 & P15 & Presión total en el ducto de bypass \\ \hline
    \filatabla s7 & P30 & Presión total a la salida del HPC \\ \hline
    \filatabla s8 & Nf & Velocidad física del fan \\ \hline
    \filatabla s9 & Nc & Velocidad física del núcleo (compresor) \\ \hline
    \filatabla s10 & epr & Razón de presión del motor \\ \hline
    \filatabla s11 & Ps30 & Presión estática a la salida del HPC \\ \hline
    \filatabla s12 & phi & Razón entre el flujo de combustible y Ps30 \\ \hline
    \filatabla s13 & NRf & Velocidad corregida del fan \\ \hline
    \filatabla s14 & NRc & Velocidad corregida del núcleo \\ \hline
    \filatabla s15 & BPR & Razón de bypass \\ \hline
    \filatabla s16 & farB & Razón combustible/aire del quemador \\ \hline
    \filatabla s17 & htBleed & Entalpía de sangrado \\ \hline
    \filatabla s18 & Nf\_dmd & Velocidad del fan demandada \\ \hline
    \filatabla s19 & PCNfR\_dmd & Velocidad corregida del fan demandada \\ \hline
    \filatabla s20 & W31 & Sangrado de enfriamiento de la turbina de alta (HPT) \\ \hline
    \filatabla s21 & W32 & Sangrado de enfriamiento de la turbina de baja (LPT) \\ \hline
\end{tabularx}
\end{table}

\subsection{Tratamiento automático de los datos}

Al entrenar, el sistema realiza de forma automática las siguientes operaciones
(Módulos~1 y~2):

\begin{enumerate}[leftmargin=*, label=\arabic*.]
    \item \textbf{Partición por motor:} el 20\,\% de los motores de entrenamiento
          se reserva para validación interna. La partición se hace por motor, no
          por ventana, para evitar que ventanas casi idénticas de un mismo motor
          queden en ambos conjuntos.
    \item \textbf{Eliminación de señales constantes:} se descarta un sensor si su
          desviación estándar dentro de cada régimen de operación es menor que
          0.01 en todos los regímenes. La lista de sensores conservados y
          descartados se guarda en \ruta{metrics.json}.
    \item \textbf{Normalización por régimen de operación:} escalamiento mínimo--máximo
          con estadísticas calculadas \emph{solo} con los motores de entrenamiento.
          En FD002 y FD004 los seis regímenes se identifican con K-Means sobre
          los ajustes operacionales.
    \item \textbf{Etiqueta de RUL lineal por tramos:} $\text{RUL}(t)=t_{\text{final}}-t$,
          limitada a un tope (por defecto 125 ciclos)~\cite{heimes2008rnn}.
    \item \textbf{Ventanas deslizantes:} ventanas de longitud fija (por defecto 30
          ciclos, paso 1); la etiqueta es la RUL del último ciclo de la ventana.
          En prueba se usa la última ventana de cada motor; si la trayectoria es
          más corta que la ventana, se rellena replicando la primera lectura.
\end{enumerate}

% ============================================================
\section{Entrenamiento y evaluación del modelo}
\label{sec:entrenamiento}

\subsection{Entrenar un subconjunto}

\begin{lstlisting}
make train SUBSET=FD001 WINDOW=30
\end{lstlisting}

Para entrenar los cuatro subconjuntos con las ventanas predeterminadas
(30, 20, 30 y 20 ciclos para FD001 a FD004, respectivamente):

\begin{lstlisting}
make train-all
\end{lstlisting}

El proceso ejecuta, en orden: carga y limpieza, construcción de ventanas,
entrenamiento con parada temprana, evaluación en el conjunto de prueba, y
guardado de modelo, métricas y gráficas.

\subsection{Arquitectura e hiperparámetros por defecto}

El modelo es una red LSTM apilada seguida de capas densas y una salida lineal que
entrega la RUL en ciclos~\cite{zheng2017lstm,hochreiter1997lstm}. El
entrenamiento usa el optimizador Adam~\cite{kingma2015adam} con tasa de
aprendizaje inicial $10^{-3}$, reducción adaptativa de la tasa al detectarse
estancamiento de la pérdida de validación (\textit{ReduceLROnPlateau}) y parada
temprana. Para prevenir el sobreajuste se emplea abandono
(\textit{dropout})~\cite{srivastava2014dropout}.

\begin{table}[H]
\centering
\fontsize{9}{11}\selectfont
\caption{Parámetros de \cmd{python -m src.train}.}
\label{tab:parametros}
\renewcommand{\arraystretch}{1.5}
\begin{tabularx}{\textwidth}{|>{\raggedright\arraybackslash}p{3.7cm}|>{\raggedright\arraybackslash}p{2.6cm}|X|}
    \hline
    \encabezadotabla \textbf{Parámetro} & \textbf{Valor por defecto} & \textbf{Descripción} \\
    \hline
    \filatabla \cmd{--subset} & FD001 & Subconjunto: FD001, FD002, FD003 o FD004 \\ \hline
    \filatabla \cmd{--window} & 30 & Longitud de la ventana temporal (ciclos) \\ \hline
    \filatabla \cmd{--rul-cap} & 125 & Tope de la RUL lineal por tramos \\ \hline
    \filatabla \cmd{--val-fraction} & 0.2 & Fracción de motores para validación interna \\ \hline
    \filatabla \cmd{--lstm-units} & 64 32 & Celdas de cada capa LSTM (una cifra por capa) \\ \hline
    \filatabla \cmd{--dense-units} & 32 16 & Neuronas de cada capa densa \\ \hline
    \filatabla \cmd{--dropout} & 0.2 & Tasa de abandono \\ \hline
    \filatabla \cmd{--learning-rate} & 0.001 & Tasa de aprendizaje inicial de Adam \\ \hline
    \filatabla \cmd{--epochs} & 100 & Máximo de épocas (la parada temprana puede detener antes) \\ \hline
    \filatabla \cmd{--batch-size} & 256 & Tamaño de lote \\ \hline
    \filatabla \cmd{--seed} & 42 & Semilla de reproducibilidad \\ \hline
    \filatabla \cmd{--raw-dir} & \ruta{data/raw} & Carpeta con los archivos de C-MAPSS \\ \hline
    \filatabla \cmd{--out-dir} & \ruta{models} & Carpeta de salida de los artefactos \\ \hline
    \filatabla \cmd{--verbose} & 1 & 0 = silencioso; 1 = muestra el progreso por época \\ \hline
\end{tabularx}
\end{table}

\noindent Ejemplo con hiperparámetros personalizados:

\begin{lstlisting}
python -m src.train --subset FD002 --window 20 --lstm-units 128 64 --dropout 0.3 --epochs 150
\end{lstlisting}

\textbf{Nota.} Los valores por defecto son un punto de partida, no los valores
finales; deben refinarse con la búsqueda en cuadrícula y contrastarse con lo
reportado en la literatura~\cite{zheng2017lstm}. La ubicación de datos y modelos
también puede cambiarse con las variables de entorno \cmd{RUL_DATA_DIR} y
\cmd{RUL_MODELS_DIR}.

\subsection{Búsqueda en cuadrícula de hiperparámetros}

\begin{lstlisting}
make grid SUBSET=FD001
\end{lstlisting}

Explora combinaciones de tamaño de ventana, celdas por capa LSTM y tasa de
abandono. \textbf{La selección se basa únicamente en la pérdida de la validación
interna}: el conjunto de prueba no interviene, para no sobreajustar la selección
a la evaluación final. Por defecto se evalúan 4 ventanas $\times$ 4
configuraciones LSTM $\times$ 3 tasas de abandono = 48 entrenamientos, lo que
puede tardar horas; para una exploración más breve:

\begin{lstlisting}
python -m src.grid_search --subset FD001 --windows 20 30 --lstm-options 32,32 64,32 --dropouts 0.1 0.2 --epochs 40
\end{lstlisting}

El resultado se guarda en \ruta{models/grid_search/grid_FD001.csv}, ordenado de
menor a mayor \cmd{val_rmse}; la primera fila es la mejor configuración. Con ella
se vuelve a ejecutar \cmd{src.train} para el entrenamiento final.

\subsection{Artefactos generados}

Cada entrenamiento crea la carpeta \ruta{models/<SUBCONJUNTO>/}:

\begin{table}[H]
\centering
\fontsize{9}{11}\selectfont
\caption{Archivos generados por el entrenamiento.}
\label{tab:artefactos}
\renewcommand{\arraystretch}{1.5}
\begin{tabularx}{\textwidth}{|>{\raggedright\arraybackslash}p{4.6cm}|X|}
    \hline
    \encabezadotabla \textbf{Archivo} & \textbf{Contenido} \\
    \hline
    \filatabla \ruta{model.keras} & Modelo LSTM entrenado (devuelve RUL en ciclos) \\ \hline
    \filatabla \ruta{preprocessor.joblib} & Sensores seleccionados, regímenes y parámetros de normalización \\ \hline
    \filatabla \ruta{config.json} & Hiperparámetros del experimento \\ \hline
    \filatabla \ruta{metrics.json} & Métricas de validación y de prueba; sensores usados y descartados \\ \hline
    \filatabla \ruta{history.csv} & Pérdida y métricas por época \\ \hline
    \filatabla \ruta{test_predictions.csv} & RUL real y estimada de cada motor de prueba \\ \hline
    \filatabla \ruta{figures/pred_vs_real.png} & RUL real vs. estimada (motores ordenados por RUL real) y dispersión \\ \hline
    \filatabla \ruta{figures/error_dist.png} & Histograma del error de estimación \\ \hline
    \filatabla \ruta{figures/training_history.png} & Curvas de pérdida y de la tasa de aprendizaje \\ \hline
\end{tabularx}
\end{table}

Para reevaluar un modelo ya entrenado sobre el conjunto de prueba:

\begin{lstlisting}
make evaluate SUBSET=FD001
\end{lstlisting}

\subsection{Interpretación de las métricas}

Sea $d_i=\widehat{\text{RUL}}_i-\text{RUL}_i$ el error del motor $i$ (RUL estimada
menos RUL real) y $N$ el número de motores de prueba. Se calculan las
siguientes métricas~\cite{saxena2008damage}:

\begin{align}
\text{RMSE} &= \sqrt{\frac{1}{N}\sum_{i=1}^{N} d_i^{2}}, &
\text{MAE} &= \frac{1}{N}\sum_{i=1}^{N} |d_i|, \\[4pt]
\text{Score} &= \sum_{i=1}^{N} s_i, &
s_i &=
\begin{cases}
e^{-d_i/13}-1, & d_i<0 \quad \text{(estimación temprana)}\\[2pt]
e^{\,d_i/10}-1, & d_i\geq 0 \quad \text{(estimación tardía)}
\end{cases}
\end{align}

\begin{itemize}
    \item \textbf{RMSE} y \textbf{MAE} se expresan en ciclos; menor es mejor. El
          RMSE pondera más los errores grandes.
    \item \textbf{Score} es una función de puntuación asimétrica; menor es
          mejor. Penaliza más las estimaciones \textbf{tardías}
          ($d_i>0$, la RUL se \emph{sobreestima}), porque implican operar el
          motor más allá de su vida segura, que las tempranas ($d_i<0$).
    \item \cmd{pct_estimacion_tardia} indica el porcentaje de motores con
          $d_i>0$; un valor muy superior a 50\,\% señala un modelo sesgado
          hacia el optimismo.
\end{itemize}

En \ruta{metrics.json} se reportan tres bloques: \cmd{validation_windows}
(todas las ventanas de los motores de validación), \cmd{test} (última ventana de
cada motor de prueba contra la RUL real \emph{sin tope}) y
\cmd{test_capped_labels} (igual, pero con la RUL real limitada al tope). Al
comparar con otros trabajos, indique cuál convención utiliza, porque la
literatura no es uniforme.

% ============================================================
\section{Uso del tablero de visualización}
\label{sec:tablero}

\subsection{Iniciar el tablero}

Requisito: haber entrenado al menos un subconjunto (Sección~\ref{sec:entrenamiento})
y contar con los datos en \ruta{data/raw/}.

\begin{lstlisting}
make dashboard
\end{lstlisting}

El navegador se abre en \cmd{http://localhost:8501}. Si no se abre
automáticamente, copie esa dirección en el navegador. Para usar otro puerto:
\cmd{streamlit run dashboard/app.py --server.port 8502}. Para detener el tablero,
presione \cmd{Ctrl+C} en la terminal.

Si no hay modelos entrenados, el tablero lo indica con un mensaje y el comando
que debe ejecutarse (\cmd{make train SUBSET=FD001}).

\subsection{Distribución de la pantalla}

\begin{figure}[H]
\centering
\begin{tikzpicture}[x=1cm, y=1cm,
    marco/.style={draw=gray!60, rounded corners=3pt},
    num/.style={circle, draw=red!60!black, fill=red!10, inner sep=1.2pt,
                font=\fontsize{8}{9}\selectfont\bfseries, text=red!50!black}
  ]
  % Marco general
  \draw[marco, fill=white] (0,0) rectangle (15,10);
  % Barra lateral
  \draw[marco, fill=gray!14] (0,0) rectangle (3.4,10);
  \node[font=\fontsize{8}{9}\selectfont\bfseries, anchor=north west] at (0.2,9.8) {Configuración};
  \node[font=\fontsize{7}{8}\selectfont, anchor=north west, align=left] at (0.2,9.2)
        {Subconjunto C-MAPSS\\[2pt]\fbox{\strut FD001\hspace{0.6cm}\ding{116}}};
  \node[font=\fontsize{7}{8}\selectfont, anchor=north west, align=left] at (0.2,7.9)
        {Umbrales de alerta\\(ciclos)};
  \draw[gray!70, thick] (0.3,6.6) -- (3.1,6.6); \fill[red!60] (1.0,6.6) circle (0.09);
  \node[font=\fontsize{6.5}{7}\selectfont, anchor=west] at (0.3,6.9) {Crítico};
  \draw[gray!70, thick] (0.3,5.7) -- (3.1,5.7); \fill[orange!80] (1.9,5.7) circle (0.09);
  \node[font=\fontsize{6.5}{7}\selectfont, anchor=west] at (0.3,6.0) {Precaución};
  % Pestañas
  \foreach \x/\t in {3.7/Motor individual, 6.3/Flota (prueba), 8.9/Modelo}{
    \draw[marco, fill=colorPropuesta] (\x,9.1) rectangle (\x+2.4,9.8);
    \node[font=\fontsize{7.5}{8}\selectfont] at (\x+1.2,9.45) {\t};
  }
  % Banner de alerta
  \draw[marco, fill=semCritico] (3.7,8.0) rectangle (14.8,8.8);
  \node[font=\fontsize{7.5}{8}\selectfont] at (9.25,8.4) {ALERTA DE FALLA — CRÍTICO: RUL estimada = 24 ciclos};
  % Indicador tipo velocímetro (semicírculo: 0 ciclos a la izquierda, 125 a la derecha)
  \draw[marco] (3.7,5.0) rectangle (7.7,7.8);
  \node[font=\fontsize{6.5}{7}\selectfont] at (5.7,7.55) {RUL estimada (último ciclo)};
  \def\cx{5.7}\def\cy{5.55}\def\rr{1.4}
  \draw[line width=7pt, semNormal]     ($(\cx,\cy)+(93.6:\rr)$)  arc[start angle=93.6,  end angle=0,    radius=\rr];
  \draw[line width=7pt, semPrecaucion] ($(\cx,\cy)+(136.8:\rr)$) arc[start angle=136.8, end angle=93.6, radius=\rr];
  \draw[line width=7pt, semCritico]    ($(\cx,\cy)+(180:\rr)$)   arc[start angle=180,   end angle=136.8, radius=\rr];
  \draw[thick, red!60!black, -{Stealth[length=2mm]}] (\cx,\cy) -- ++(145.4:1.15);
  \fill[red!60!black] (\cx,\cy) circle (0.07);
  \node[font=\fontsize{9}{10}\selectfont\bfseries] at (5.7,5.2) {24 ciclos};
  % Métricas
  \foreach \x/\t/\v in {7.9/Ciclos observados/160, 10.4/RUL estimada/24 ciclos, 12.9/RUL real (ref.)/30 ciclos}{
    \draw[marco] (\x,5.0) rectangle (\x+1.9,7.8);
    \node[font=\fontsize{6.5}{7}\selectfont] at (\x+0.95,7.2) {\t};
    \node[font=\fontsize{8}{9}\selectfont\bfseries] at (\x+0.95,6.2) {\v};
  }
  % Curva de RUL
  \draw[marco] (3.7,2.4) rectangle (14.8,4.8);
  \fill[semCritico, opacity=0.6] (3.9,2.55) rectangle (14.6,2.95);
  \fill[semPrecaucion, opacity=0.6] (3.9,2.95) rectangle (14.6,3.5);
  \draw[blue!70, dashed, thick] (3.9,4.1) -- (8.0,4.1) -- (14.6,2.7);
  \draw[red!70, thick] (3.9,4.1) -- (7.8,4.05) -- (9.5,3.75) -- (11.0,3.4) -- (12.8,3.0) -- (14.6,2.75);
  \node[font=\fontsize{6.5}{7}\selectfont, anchor=north west] at (3.9,4.75) {Curva de RUL: estimada (rojo) vs. real (azul discontinuo)};
  % Sensores
  \draw[marco] (3.7,0.2) rectangle (14.8,2.2);
  \draw[gray!70, thick] (3.9,1.8) -- (6.0,1.5) -- (8.0,1.3) -- (10.0,0.9) -- (12.0,0.7) -- (14.6,0.4);
  \node[font=\fontsize{6.5}{7}\selectfont, anchor=north west] at (3.9,2.15) {Evolución temporal de sensores};
  % Numeración
  \node[num] at (3.4,9.7) {1};
  \node[num] at (6.2,9.6) {2};
  \node[num] at (14.6,8.4) {3};
  \node[num] at (7.5,7.6) {4};
  \node[num] at (14.6,7.6) {5};
  \node[num] at (14.6,4.6) {6};
  \node[num] at (14.6,2.0) {7};
\end{tikzpicture}
\caption{Esquema de la pestaña «Motor individual». Los números se explican en
la Tabla~\ref{tab:elementos}. La figura es ilustrativa; los valores mostrados son
de ejemplo.}
\label{fig:tablero}
\end{figure}

\IfFileExists{figuras/tablero_motor.png}{%
\begin{figure}[H]
\centering
\includegraphics[width=\textwidth]{figuras/tablero_motor.png}
\caption{Captura de pantalla de la pestaña «Motor individual».}
\end{figure}}{}

\begin{table}[H]
\centering
\fontsize{9}{11}\selectfont
\caption{Elementos de la pestaña «Motor individual» (Figura~\ref{fig:tablero}).}
\label{tab:elementos}
\renewcommand{\arraystretch}{1.5}
\begin{tabularx}{\textwidth}{|c|>{\raggedright\arraybackslash}p{3.6cm}|X|}
    \hline
    \encabezadotabla \textbf{N.º} & \textbf{Elemento} & \textbf{Función} \\
    \hline
    \filatabla 1 & Barra lateral & Selecciona el subconjunto C-MAPSS (solo aparecen los entrenados) y ajusta los umbrales de alerta \\ \hline
    \filatabla 2 & Pestañas & «Motor individual», «Flota (prueba)» y «Modelo» \\ \hline
    \filatabla 3 & Banner de alerta & Estado del motor: NORMAL (verde), PRECAUCIÓN (amarillo) o CRÍTICO (rojo) \\ \hline
    \filatabla 4 & Indicador tipo velocímetro & RUL estimada en el último ciclo observado, con las zonas de color de los umbrales \\ \hline
    \filatabla 5 & Indicadores numéricos & Ciclos observados, RUL estimada y RUL real de referencia con su error (estimada $-$ real) \\ \hline
    \filatabla 6 & Curva de RUL & RUL estimada y RUL real en cada ciclo; bandas de color de los umbrales; zona gris de ventana incompleta \\ \hline
    \filatabla 7 & Sensores & Evolución temporal de los sensores elegidos, en valores originales o normalizados \\ \hline
\end{tabularx}
\end{table}

\subsection{Procedimiento de uso}

\begin{enumerate}[leftmargin=*, label=\arabic*.]
    \item \textbf{Elegir el subconjunto.} En la barra lateral, seleccione FD001,
          FD002, FD003 o FD004. Solo se listan los subconjuntos con un modelo
          entrenado en \ruta{models/}.
    \item \textbf{Ajustar los umbrales (opcional).} Con los deslizadores defina
          los ciclos restantes a partir de los cuales el estado pasa a
          PRECAUCIÓN y a CRÍTICO (Sección~\ref{sec:alertas}).
    \item \textbf{Seleccionar un motor} en la pestaña «Motor individual». El
          tablero calcula la RUL estimada en cada ciclo y muestra el estado
          actual.
    \item \textbf{Leer la alerta y el indicador} (elementos 3 y 4 de la
          Figura~\ref{fig:tablero}).
    \item \textbf{Analizar la curva de RUL.} Compare la estimación con la RUL
          real y observe la tendencia.
    \item \textbf{Explorar los sensores.} Elija uno o varios sensores en la lista
          (se muestra el símbolo y la descripción). Active «Mostrar valores
          normalizados» para compararlos en una escala común por régimen.
    \item \textbf{Revisar la flota.} En «Flota (prueba)» se listan todos los
          motores ordenados por RUL estimada ascendente, de modo que los de
          intervención más urgente aparecen primero; el gráfico de barras y la
          tabla se actualizan con los umbrales elegidos.
    \item \textbf{Consultar el modelo.} En «Modelo» se muestran RMSE, MAE y Score
          de prueba, los sensores utilizados y descartados, los hiperparámetros
          y la gráfica de RUL real vs. estimada.
\end{enumerate}

\IfFileExists{figuras/tablero_flota.png}{%
\begin{figure}[H]
\centering
\includegraphics[width=\textwidth]{figuras/tablero_flota.png}
\caption{Captura de pantalla de la pestaña «Flota (prueba)».}
\end{figure}}{}

\subsection{Sistema de alertas}
\label{sec:alertas}

El estado se determina comparando la RUL estimada del último ciclo con dos
umbrales configurables (Figura~\ref{fig:semaforo} y Tabla~\ref{tab:alertas}).

\begin{figure}[H]
\centering
\begin{tikzpicture}[x=0.115cm, y=1cm]
  \fill[semCritico]    (0,0)   rectangle (30,0.8);
  \fill[semPrecaucion] (30,0)  rectangle (60,0.8);
  \fill[semNormal]     (60,0)  rectangle (125,0.8);
  \draw[gray!70] (0,0) rectangle (125,0.8);
  \draw[thick] (30,-0.15) -- (30,0.95);
  \draw[thick] (60,-0.15) -- (60,0.95);
  \node[font=\fontsize{8}{9}\selectfont\bfseries] at (15,0.4) {CRÍTICO};
  \node[font=\fontsize{8}{9}\selectfont\bfseries] at (45,0.4) {PRECAUCIÓN};
  \node[font=\fontsize{8}{9}\selectfont\bfseries] at (92.5,0.4) {NORMAL};
  \foreach \v/\l in {0/0, 30/30, 60/60, 125/125}{
    \draw (\v,0) -- (\v,-0.15); \node[font=\fontsize{7}{8}\selectfont, below] at (\v,-0.15) {\l};
  }
  \node[font=\fontsize{7}{8}\selectfont, above, align=center] at (30,0.95) {umbral crítico\\(por defecto 30)};
  \node[font=\fontsize{7}{8}\selectfont, above, align=center] at (60,0.95) {umbral de precaución\\(por defecto 60)};
  \node[font=\fontsize{7}{8}\selectfont, below] at (62.5,-0.75) {RUL estimada (ciclos restantes); el máximo mostrado es el tope de RUL (125)};
\end{tikzpicture}
\caption{Zonas de alerta según la RUL estimada, con los umbrales por defecto.}
\label{fig:semaforo}
\end{figure}

\begin{table}[H]
\centering
\fontsize{9}{11}\selectfont
\caption{Estados de alerta del tablero.}
\label{tab:alertas}
\renewcommand{\arraystretch}{1.5}
\begin{tabularx}{\textwidth}{|>{\raggedright\arraybackslash}p{2.6cm}|>{\raggedright\arraybackslash}p{4.4cm}|X|}
    \hline
    \encabezadotabla \textbf{Estado} & \textbf{Condición} & \textbf{Acción sugerida} \\
    \hline
    \cellcolor{semNormal} NORMAL & RUL estimada $>$ umbral de precaución & Operación normal; seguimiento rutinario \\ \hline
    \cellcolor{semPrecaucion} PRECAUCIÓN & Umbral crítico $<$ RUL estimada $\leq$ umbral de precaución & Planificar el mantenimiento y reservar recursos \\ \hline
    \cellcolor{semCritico} CRÍTICO & RUL estimada $\leq$ umbral crítico & Programar la intervención en la próxima escala planificada \\ \hline
\end{tabularx}
\end{table}

\textbf{Los umbrales por defecto (60 y 30 ciclos) no provienen de la
literatura:} son parámetros de política de mantenimiento que debe definir el
departamento usuario según sus costos, tiempos de reposición y escalas
disponibles.

\subsection{Cómo interpretar la curva de RUL}

\begin{itemize}
    \item Al inicio de la vida del motor, tanto la RUL real como la estimada se
          muestran limitadas al tope (125 ciclos): se asume que mientras el motor
          está sano la degradación es despreciable~\cite{heimes2008rnn}.
    \item La zona gris «ventana incompleta» corresponde a los primeros ciclos,
          en los que la ventana de entrada se completa replicando la primera
          lectura; las estimaciones ahí son \textbf{menos confiables}.
    \item Es más informativa la \textbf{tendencia descendente sostenida} que un
          valor puntual: una estimación aislada puede oscilar por ruido.
    \item La RUL real de referencia solo existe para los motores del conjunto de
          prueba de C-MAPSS; con datos reales de operación no estaría disponible.
    \item Un error positivo (estimada $>$ real) es una estimación \textbf{tardía}
          y es el caso de mayor riesgo operativo.
\end{itemize}

\subsection{Actualizar el tablero después de reentrenar}

El tablero guarda en memoria el modelo de cada subconjunto. Después de volver a
entrenar, deténgalo (\cmd{Ctrl+C}) y vuelva a iniciarlo, o use el menú
\textit{Clear cache} de Streamlit y recargue la página.

% ============================================================
\section{Estructura del repositorio}
\label{sec:estructura}

\begin{lstlisting}
rul-turbofan-tt/
|-- data/
|   |-- raw/                  # C-MAPSS original (no versionado)
|   `-- processed/            # artefactos intermedios (no versionado)
|-- docs/
|   `-- manual_usuario/       # este manual (.tex, .pdf, figuras/)
|-- notebooks/
|   |-- 01_EDA.ipynb          # analisis exploratorio
|   `-- 02_experimentacion.ipynb
|-- src/                      # modulos del prototipo
|   |-- config.py  ingestion.py  preprocessing.py
|   |-- model.py  train.py  grid_search.py
|   `-- evaluation.py  inference.py
|-- models/<FD00X>/           # modelo, preprocesador, metricas y figuras
|-- dashboard/app.py          # tablero Streamlit
|-- tests/                    # pruebas con pytest (24)
|-- Makefile  README.md  requirements*.txt  pytest.ini  ruff.toml
`-- .gitignore
\end{lstlisting}

\subsection{Compilar este manual}

\begin{lstlisting}
make manual
\end{lstlisting}

Requiere una distribución \LaTeX{} con \cmd{latexmk} y los paquetes
\textit{tikz}, \textit{babel}, \textit{diagbox}, \textit{listings} y
\textit{makecell} (incluidos en TeX Live y MiKTeX completos). El resultado es
\ruta{docs/manual_usuario/manual_usuario.pdf}. Para incluir capturas reales,
guárdelas como \ruta{figuras/tablero_motor.png} y \ruta{figuras/tablero_flota.png}
dentro de la carpeta del manual; se insertarán automáticamente al compilar.

% ============================================================
\section{Solución de problemas}
\label{sec:problemas}

\begin{table}[H]
\centering
\fontsize{8.5}{10.5}\selectfont
\caption{Problemas frecuentes y su solución.}
\label{tab:problemas}
\renewcommand{\arraystretch}{1.5}
\begin{tabularx}{\textwidth}{|>{\raggedright\arraybackslash}p{4.4cm}|>{\raggedright\arraybackslash}p{4.4cm}|X|}
    \hline
    \encabezadotabla \textbf{Síntoma} & \textbf{Causa probable} & \textbf{Solución} \\
    \hline
    \filatabla \cmd{FileNotFoundError: Faltan [...] en data/raw} & Los datos no se han descargado o están en otra carpeta & Ejecute \cmd{make data} o copie los archivos manualmente (Sección~\ref{sec:datos}); verifique \cmd{RUL_DATA_DIR} \\ \hline
    \filatabla La descarga con \cmd{make data} falla & La URL de la NASA cambió o no hay conexión & Descargue el conjunto manualmente o use \cmd{--url} con la dirección vigente \\ \hline
    \filatabla \cmd{make: command not found} (Windows) & \cmd{make} no está instalado & Use los comandos directos de la Tabla~\ref{tab:equivalencias} \\ \hline
    \filatabla \cmd{ModuleNotFoundError: No module named 'src'} & El comando se ejecutó fuera de la raíz del repositorio & Ubíquese en la carpeta que contiene \ruta{Makefile} y use \cmd{python -m src.train} \\ \hline
    \filatabla Error al importar TensorFlow o Keras & Versión de Python no compatible o instalación incompleta & Use Python 3.10--3.12 en un entorno virtual limpio y reinstale con \cmd{make setup} \\ \hline
    \filatabla El tablero muestra «No se encontraron modelos entrenados» & No se ha entrenado ningún subconjunto & Ejecute \cmd{make train SUBSET=FD001} y recargue la página \\ \hline
    \filatabla El puerto 8501 está ocupado & Hay otra instancia de Streamlit en ejecución & Use \cmd{--server.port 8502} o cierre la instancia anterior \\ \hline
    \filatabla El tablero no refleja un modelo recién entrenado & Caché en memoria del tablero & Reinicie el tablero o use \textit{Clear cache} \\ \hline
    \filatabla Memoria insuficiente al entrenar & Lote o modelo demasiado grandes para el equipo & Reduzca \cmd{--batch-size} (p.~ej., 64) o las celdas de \cmd{--lstm-units} \\ \hline
    \filatabla RMSE muy alto o la pérdida no disminuye & Hiperparámetros poco adecuados o pocas épocas & Revise \ruta{figures/training_history.png}; ejecute \cmd{make grid}; aumente \cmd{--epochs} \\ \hline
    \filatabla La búsqueda en cuadrícula tarda demasiado & Se evalúan 48 combinaciones por defecto & Reduzca las listas de \cmd{--windows}, \cmd{--lstm-options} y \cmd{--dropouts}, y las épocas \\ \hline
    \filatabla Resultados distintos entre ejecuciones & Semilla distinta o diferencias de hardware & Fije \cmd{--seed}; pequeñas diferencias numéricas entre equipos son esperables \\ \hline
\end{tabularx}
\end{table}

% ============================================================
\section{Limitaciones y uso responsable}
\label{sec:limitaciones}

\begin{itemize}
    \item \textbf{Datos simulados.} Los resultados corresponden a C-MAPSS y no
          garantizan el desempeño sobre flotas reales ni sobre un modelo de
          aeronave específico.
    \item \textbf{Prototipo académico.} El sistema es una herramienta de apoyo a
          la decisión y de demostración; no está certificado ni sustituye las
          inspecciones y los programas de mantenimiento aprobados.
    \item \textbf{Umbrales de alerta no calibrados} con costos operativos reales
          (Sección~\ref{sec:alertas}).
    \item \textbf{Incertidumbre.} La estimación es puntual; el sistema no reporta
          intervalos de confianza.
    \item \textbf{Fidelidad de las comparaciones.} Los resultados deben
          compararse con otros trabajos solo si comparten el protocolo de
          evaluación (tope de RUL, tamaño de ventana, uso del conjunto de prueba).
\end{itemize}

% ============================================================
\section{Glosario}

\begin{table}[H]
\centering
\fontsize{9}{11}\selectfont
\caption{Términos utilizados en este manual.}
\label{tab:glosario}
\renewcommand{\arraystretch}{1.45}
\begin{tabularx}{\textwidth}{|>{\raggedright\arraybackslash}p{3.8cm}|X|}
    \hline
    \encabezadotabla \textbf{Término} & \textbf{Definición} \\
    \hline
    \filatabla C-MAPSS & \textit{Commercial Modular Aero-Propulsion System Simulation}: simulador y conjuntos de datos de degradación de motores turbofan de la NASA \\ \hline
    \filatabla Ciclo & Una unidad de operación del motor (por ejemplo, un vuelo) en el que se registran las lecturas de los sensores \\ \hline
    \filatabla RUL & \textit{Remaining Useful Life}: ciclos de operación que restan antes de la falla \\ \hline
    \filatabla Tope de RUL & Valor máximo de la RUL usada como etiqueta (RUL lineal por tramos); por defecto 125 ciclos \\ \hline
    \filatabla Régimen de operación & Combinación de condiciones de vuelo (ajustes operacionales); FD002 y FD004 tienen seis \\ \hline
    \filatabla Ventana temporal & Secuencia de los últimos $W$ ciclos de lecturas que recibe la red para estimar la RUL \\ \hline
    \filatabla LSTM & \textit{Long Short-Term Memory}: red neuronal recurrente capaz de aprender dependencias de largo plazo~\cite{hochreiter1997lstm} \\ \hline
    \filatabla Dropout (abandono) & Técnica de regularización que desactiva neuronas al azar durante el entrenamiento para reducir el sobreajuste~\cite{srivastava2014dropout} \\ \hline
    \filatabla Validación interna & 20\,\% de los motores de entrenamiento reservados para seleccionar hiperparámetros y detener el entrenamiento \\ \hline
    \filatabla Estimación tardía / temprana & Error positivo ($\widehat{\text{RUL}}>\text{RUL}$) o negativo, respectivamente \\ \hline
    \filatabla RMSE, MAE, Score & Métricas de desempeño definidas en la Sección~\ref{sec:entrenamiento} \\ \hline
    \filatabla CRISP-DM & Proceso estándar de la industria para proyectos de minería de datos~\cite{wirth2000crisp} \\ \hline
\end{tabularx}
\end{table}

% ============================================================
\begin{thebibliography}{9}

\bibitem{frederick2007cmapss}
Frederick, D., de Castro, J. y Litt, J. (2007).
\textit{User's Guide for the Commercial Modular Aero-Propulsion System
Simulation (C-MAPSS)}.
NASA Technical Memorandum NASA/TM--2007-215026. Glenn Research Center,
Cleveland, OH: NASA.

\bibitem{heimes2008rnn}
Heimes, F.~O. (2008).
Recurrent neural networks for remaining useful life estimation.
En \textit{2008 International Conference on Prognostics and Health
Management (PHM08)} (pp.~1--6). Denver, CO: IEEE.

\bibitem{hochreiter1997lstm}
Hochreiter, S. y Schmidhuber, J. (1997).
Long short-term memory.
\textit{Neural Computation}, 9(8), 1735--1780.
\url{https://doi.org/10.1162/neco.1997.9.8.1735}

\bibitem{kingma2015adam}
Kingma, D.~P. y Ba, J. (2015).
Adam: A method for stochastic optimization.
En \textit{3rd International Conference on Learning Representations
(ICLR 2015)}. San Diego, CA.
\url{https://arxiv.org/abs/1412.6980}

\bibitem{saxena2008damage}
Saxena, A., Goebel, K., Simon, D. y Eklund, N. (2008).
Damage propagation modeling for aircraft engine run-to-failure simulation.
En \textit{Proceedings of the 1st International Conference on Prognostics
and Health Management (PHM08)}. Denver, CO: IEEE.
\url{https://doi.org/10.1109/PHM.2008.4711414}

\bibitem{srivastava2014dropout}
Srivastava, N., Hinton, G., Krizhevsky, A., Sutskever, I. y
Salakhutdinov, R. (2014).
Dropout: A simple way to prevent neural networks from overfitting.
\textit{Journal of Machine Learning Research}, 15(56), 1929--1958.

\bibitem{wirth2000crisp}
Wirth, R. y Hipp, J. (2000).
CRISP-DM: Towards a standard process model for data mining.
En \textit{Proceedings of the 4th International Conference on the
Practical Applications of Knowledge Discovery and Data Mining},
Manchester, 29--39.

\bibitem{zheng2017lstm}
Zheng, S., Ristovski, K., Farahat, A. y Gupta, C. (2017).
Long short-term memory network for remaining useful life estimation.
En \textit{2017 IEEE International Conference on Prognostics and Health
Management (ICPHM)} (pp.~88--95). Dallas, TX: IEEE.

\end{thebibliography}

\end{document}